% Nouveau EGA V (cote 25): archive cote 25, site page 159 (PDF page 160).
% English translation of the checked French annotated master.
% AI-translated by GPT-6 Astra (OpenAI, ChatGPT Pro mode).
\scanpage{25/159}
\mspage{\hand{2}}

singular zero" if one insists on terminology
close to the terminology accepted in char.0; likewise, one may
speak of a "nondegenerate quadratic singularity" of a loc
noeth prescheme, resp. a "geo\unsure{m}\unsure{e}trically nondegenerate quadratic singularity" when it is
loc of finite type over a field k. On the other hand, the terminology
"ordinary singular zero", or better "ordinary quadratic zero",
and likewise "ordinary quadratic singularity", "geometr\unsure{i}cally ordinary quadratic si\unsure{n}gularity",
must be understood (in accordance with usage)
as corresponding to the notion of a smooth quadratic form. Moreover,
it may be appropriate to replace the term "ordinary" by "elementary",
and to extend this terminology to singularities that are not necessarily
quadratic but have arbitrary multiplicity. You will tell me your impression
on this point. --- It does indeed seem, on the other hand, that the text SUP is
formally correct; in particular, in\hand{/}prop.7 the notion that actually enters
is that of a supersingular zero. Nevertheless, it would be appropriate,
at least in a remark, also to introduce \var{\xout{$V(\hand{\phi})^{\unsure{\text{sing él}}}$}}{}{typewriter strike-out}{ega25_25_159-v0002} the subprescheme
$V(\hand{\phi})^{\text{ult sing}}$ of $V(\hand{\phi})^{\text{sup sing}}$, corresponding to the consideration
of \unsure{ultrasingular} zeros, i.e. \var{\xout{not}}{}{typewriter strike-out}{ega25_25_159-v0003} which are singular without being
elementary quadratic; it is described essentially by the same
procedure as $V(\hand{\phi})^{\text{sup sing}}$, by taking the prescheme of zeros of the
"corrected discriminant" introduced below, instead of the ordinary
discriminant. Thus the only corrections to be made to SUP seem to concern
terminology (and they must also apply, of course, to \var{\xout{\unsure{that}}}{}{typewriter strike-out}{ega25_25_159-v0004}\var{\inserted[typed]{the}}{the}{insertion}{ega25_25_159-v0005}
\var{\inserted[typed]{terminology}}{terminology}{insertion}{ega25_25_159-v0006}
introduced in SUP p.10\hand{/}\var{\inserted{above}}{above}{insertion}{ega25_25_159-v0007}). On the other hand, IV 23.9.2. is false as it stands
unless k has char $\hand{\neq} 2$ \uline{or} the irreducible components of X have
even dimension; in the general case, one must delete "$f_n$ satisfies
the equivalent conditions of 8.8., in particular \var{\deleted{\illegible{1}}}{}{deletion}{ega25_25_159-v0008}"; the rest of the

\mspage{\hand{159}}
\tnote{25-159-1}{The opening completes the quoted term begun on the preceding scan, with the ordinary English order of its adjectives. The source has sin for si and unaccented recue; these do not create English spelling errors. Its canceled doubtful \unsure{ré} is an incomplete French fragment with no separate English lexical content. Loc, noeth, resp., char and prop. mean locally, noetherian, respectively, characteristic and proposition.}
\tnote{25-159-2}{The two linked doubts in geometrically retain the French candidates m/é and é/m. The later linked doubts in geometrically and singularity retain the source's i/m and n/g ambiguities; their English positions follow adjective order. Ultrasingular retains the doubt between ultra-singuliers and ultra-singulaers. The literal mathematical labels sing él, ult sing and sup sing abbreviate elementary singular, ultrasingular and supersingular. The reference IV 23.9.2. remains as written.}
